\documentclass[11pt,a4paper,parskip=half]{scrreprt}

% ---------- Core packages ----------
\usepackage[T1]{fontenc}
\usepackage[utf8]{inputenc}
\usepackage[a4paper,left=2.3cm,right=2.3cm,top=2.6cm,bottom=2.8cm,headsep=10pt]{geometry}
\usepackage[table]{xcolor}
\usepackage{colortbl}
\usepackage{booktabs}
\usepackage{array}
\usepackage{tabularx}
\usepackage{enumitem}
\usepackage{amsmath}
\usepackage{amssymb}
\usepackage{textcomp}
\usepackage{fontawesome5}
\usepackage{tikz}
\usetikzlibrary{positioning,calc,shapes.geometric,decorations.pathmorphing,backgrounds,patterns,arrows.meta}
\usepackage[most]{tcolorbox}
\usepackage[headsepline,footsepline]{scrlayer-scrpage}
\usepackage[colorlinks=true,linkcolor=eublue,urlcolor=teal,citecolor=eublue]{hyperref}

% ---------- Colour palette ----------
\definecolor{eublue}{RGB}{0,51,153}
\definecolor{eugold}{RGB}{255,183,0}
\definecolor{slate}{RGB}{68,76,86}
\definecolor{teal}{RGB}{0,121,140}
\definecolor{lightgray}{RGB}{244,246,249}
\definecolor{riskLow}{RGB}{206,232,214}
\definecolor{riskMed}{RGB}{255,229,180}
\definecolor{riskHigh}{RGB}{248,201,201}

\hypersetup{
  pdftitle={FLEXGRID -- Horizon Europe Innovation Action Proposal},
  pdfauthor={FLEXGRID Consortium}
}

% ---------- Heading styling ----------
\setkomafont{disposition}{\normalfont\sffamily\color{eublue}}
\addtokomafont{chapter}{\Large}
\addtokomafont{section}{\large}
\addtokomafont{subsection}{\normalsize\color{slate}}
\renewcommand*{\chapterformat}{\textcolor{eugold}{\Huge\thechapter}\enskip}

% ---------- Header / footer ----------
\pagestyle{scrheadings}
\clearpairofpagestyles
\automark[section]{chapter}
\ihead{\small\color{slate}\sffamily FLEXGRID\, \textbar\, Horizon Europe Innovation Action}
\ohead{\small\color{slate}\sffamily\rightmark}
\ifoot{\small\color{slate} Call: HORIZON-CL5-2026-D3-01}
\ofoot{\small\color{slate} Page \thepage}
\setkomafont{pageheadfoot}{\upshape}
\setkomafont{headsepline}{\color{eugold}}
\setkomafont{footsepline}{\color{eugold}}

% ---------- tcolorbox styles ----------
\tcbset{
  factbox/.style={colback=lightgray,colframe=eublue,boxrule=0.8pt,arc=2mm,
    left=7pt,right=7pt,top=6pt,bottom=6pt,
    fonttitle=\bfseries\sffamily\color{white},colbacktitle=eublue,coltitle=white,
    title={#1}},
  objbox/.style={colback=white,colframe=teal,boxrule=1pt,arc=1mm,
    left=8pt,right=8pt,top=6pt,bottom=6pt,
    fonttitle=\bfseries\sffamily\color{teal},coltitle=teal},
  notebox/.style={colback=eugold!12,colframe=eugold!70!black,boxrule=0.8pt,arc=1mm,
    left=7pt,right=7pt,top=5pt,bottom=5pt}
}

\setlist[itemize,1]{label=\textcolor{teal}{\small\faIcon{check}},leftmargin=1.5em}
\setlist[enumerate,1]{leftmargin=1.6em}

\newcommand{\eur}[1]{\texteuro#1}

\begin{document}

% ============================================================
% COVER PAGE
% ============================================================
\begin{titlepage}
\newgeometry{left=2cm,right=2cm,top=1.6cm,bottom=1.8cm}
\thispagestyle{empty}

\begin{tikzpicture}[remember picture,overlay]
  \fill[eublue] (current page.north west) rectangle ([yshift=-3.4cm]current page.north east);
  \foreach \x in {1,...,9}{
    \fill[eugold,opacity=0.55] ([xshift=\x*1.6cm+0.4cm,yshift=-0.55cm]current page.north west) circle (2.6pt);
  }
  \fill[teal,opacity=0.9] ([xshift=-1.2cm,yshift=1.2cm]current page.south east)
     -- ++(3.2cm,0) -- ++(0,3.2cm) -- cycle;
  \begin{scope}[xshift=2cm,yshift=-2.6cm]
    \foreach \i in {0,...,5}{
      \draw[pattern=north east lines,pattern color=eugold!35,draw=none]
        ([xshift=\i*0.55cm]current page.south west) rectangle ++(0.35cm,0.9cm);
    }
  \end{scope}
\end{tikzpicture}

\vspace*{0.3cm}
\begin{flushleft}
{\color{white}\sffamily\small HORIZON EUROPE \textbar\ CLUSTER 5: CLIMATE, ENERGY \& MOBILITY}\\[2pt]
{\color{eugold}\sffamily\footnotesize CALL HORIZON-CL5-2026-D3-01 \textbar\ INNOVATION ACTION (IA)}
\end{flushleft}

\vspace{2.4cm}

\begin{center}
{\Huge\sffamily\bfseries\color{eublue}FLEXGRID}\\[8pt]
{\large\sffamily\color{slate} Federated Learning for EXtended Grid resilience and\\ Intelligent Demand-response}\\[14pt]
\begin{tikzpicture}
  \draw[decorate,decoration={random steps,segment length=3.2mm,amplitude=0.6mm},
    ultra thick,eugold] (0,0) -- (10.2,0);
\end{tikzpicture}
\end{center}

\vspace{1.1cm}

\begin{center}
\begin{minipage}{0.86\textwidth}
\begin{tcolorbox}[factbox={\faIcon{file}\ \ Proposal Key Facts}]
\small
\renewcommand{\arraystretch}{1.35}
\begin{tabularx}{\linewidth}{@{}l X l X@{}}
\textbf{Duration:} & 36 months & \textbf{Type of action:} & Innovation Action (IA) \\
\textbf{Requested EU contribution:} & \eur{5,980,000} & \textbf{Consortium size:} & 6 beneficiaries, 5 countries \\
\textbf{Coordinator:} & Nexa Dynamics B.V. (NL) & \textbf{Start date (planned):} & 1 January 2027 \\
\textbf{TRL progression:} & TRL 4 $\rightarrow$ TRL 7 & \textbf{Panel:} & HORIZON-CL5-2026-D3 \\
\end{tabularx}
\end{tcolorbox}
\end{minipage}
\end{center}

\vspace{0.9cm}

\begin{center}
\small\sffamily\color{slate}
Coordinating person: Dr.\ Elena Marsh \textbar\ \faIcon{envelope}\ e.marsh@nexadynamics.example \textbar\ \faIcon{phone}\ +31\,20\,555\,0142
\end{center}

\vfill

\begin{center}
{\footnotesize\color{slate}\sffamily This document is a full proposal submitted under the Horizon Europe Work Programme 2026--2027.\\
Consortium: Nexa Dynamics (NL) \textbullet\ Synapse Analytics (DE) \textbullet\ Vertex Materials Institute (IE) \textbullet\ Nimbus Grid Systems (FI) \textbullet\ Apex Photonics Lab (SE) \textbullet\ Meridian Energy Cooperative (ES)}
\end{center}

\restoregeometry
\end{titlepage}

\clearpage
\thispagestyle{empty}
\tableofcontents
\clearpage

% ============================================================
% CHAPTER 1 -- EXCELLENCE
% ============================================================
\chapter{Excellence}

\section{Objectives and Ambition}

European distribution grids are absorbing rooftop solar, heat pumps and EV chargers faster than
control-room software can model them. Local congestion events on medium-voltage feeders have
risen 34\% across the five FLEXGRID demonstration regions since 2023, and today's SCADA-era
forecasting tools update on a 15-minute cycle -- too slow to prevent nuisance curtailment or
localised brownouts. FLEXGRID's ambition is to replace this cycle with a \textbf{federated,
edge-native intelligence layer} that lets distribution system operators (DSOs) forecast, explain
and act on grid stress in under 90 seconds, without centralising raw household data.

\begin{tcolorbox}[objbox={\faIcon{lightbulb}\ \ Specific, Measurable Objectives}]
\begin{enumerate}[label=\textbf{O\arabic*.}]
  \item Reduce feeder-level forecast error (RMSE) by \textbf{40\%} relative to the current
        state-of-the-art centralised model, validated on three real DSO territories.
  \item Cut raw household energy data leaving the edge device to \textbf{zero}, using federated
        learning with differential privacy ($\varepsilon \leq 3$).
  \item Demonstrate autonomous congestion mitigation across \textbf{$\geq$ 1{,}200} grid-edge
        nodes at TRL~7, sustaining a 99.5\% control-loop uptime over a 9-month field trial.
  \item Cut the cost of grid-edge sensing hardware by \textbf{35\%} through the photonic sensor
        design developed in WP4, enabling deployment at rural feeders previously judged
        uneconomical.
  \item Deliver an open reference architecture adopted by \textbf{at least 2} DSOs beyond the
        consortium within 12 months of project end.
\end{enumerate}
\end{tcolorbox}

These objectives extend directly beyond the state of the art: existing federated energy-forecasting
research (e.g.\ cluster-level load prediction pilots reported in 2024--2025) has stopped at
simulation or single-utility trials. FLEXGRID is the first project to combine (i) cross-border
federated learning across independently regulated DSOs, (ii) a purpose-built low-power photonic
sensing node that halves edge compute cost, and (iii) closed-loop autonomous mitigation rather than
forecasting alone.

\subsection{Methodology and Concept}

The technical approach separates learning from actuation, so that a mis-trained model can never
directly command grid hardware without a rule-based safety envelope. Figure~\ref{fig:concept}
summarises the architecture matured across WP2--WP4.

\begin{center}
\begin{tikzpicture}[
  node distance=8mm and 10mm,
  box/.style={draw=eublue,thick,fill=lightgray,rounded corners=2pt,align=center,
    minimum height=1.05cm,minimum width=2.9cm,font=\small\sffamily},
  edgebox/.style={box,fill=teal!12,draw=teal},
  core/.style={box,fill=eugold!18,draw=eugold!70!black},
  arrow/.style={-{Latex[length=2.2mm]},thick,color=slate}
]
\node[edgebox] (sensor) {Photonic sensor\\ grid-edge node};
\node[edgebox,right=of sensor] (local) {Local model\\ update (on-device)};
\node[core,right=of local] (agg) {Federated\\ aggregator};
\node[core,below=of agg] (safety) {Rule-based\\ safety envelope};
\node[edgebox,below=of local] (act) {Local\\ actuation};

\draw[arrow] (sensor) -- (local);
\draw[arrow] (local) -- (agg) node[midway,above,font=\tiny\sffamily] {gradients only};
\draw[arrow] (agg) -- (safety) node[midway,right,font=\tiny\sffamily] {global model};
\draw[arrow] (safety) -- (act);
\draw[arrow] (act) -- (sensor) node[midway,left,font=\tiny\sffamily] {feedback};
\end{tikzpicture}
\captionof{figure}{FLEXGRID closed-loop concept: raw data never leaves the edge node; only
model gradients are federated.}
\label{fig:concept}
\end{center}

\section{Relation to the Work Programme}

The proposal responds directly to destination \emph{``Cross-sectoral solutions for the climate
transition''} of Cluster~5, addressing the expected outcome that DSOs gain ``real-time,
privacy-preserving decision-support tools for distribution grid flexibility.'' FLEXGRID's
TRL~4$\rightarrow$7 trajectory and multi-country field validation match the call's requirement
for at least three demonstration sites in different Member States operating under different
regulatory frameworks.

% ============================================================
% CHAPTER 2 -- IMPACT
% ============================================================
\chapter{Impact}

\section{Expected Outcomes and Pathways to Impact}

\begin{table}[h]
\small
\renewcommand{\arraystretch}{1.3}
\begin{tabularx}{\linewidth}{@{}>{\sffamily\bfseries}p{2.1cm} X X@{}}
\toprule
\rowcolor{eublue}
\textcolor{white}{Horizon} & \textcolor{white}{Output (project end)} & \textcolor{white}{Outcome / Impact (2--5 yrs after)} \\
\midrule
Short-term & Validated federated forecasting model, open reference architecture, photonic
sensor design files & 2 additional DSOs pilot the architecture; 40\% forecast error reduction
independently confirmed \\
\addlinespace
Medium-term & Standardisation contribution to CENELEC TC~205 edge-interoperability profile &
Grid-edge sensor cost curve shifts, enabling rural feeder monitoring previously unfunded \\
\addlinespace
Long-term & Consortium spin-out offering FLEXGRID-as-a-service to mid-size DSOs & Measurable
reduction in curtailment-related renewable energy loss across early-adopter territories \\
\bottomrule
\end{tabularx}
\caption{Pathways from project outputs to sector-level impact.}
\end{table}

\subsection{Contribution to EU Policy Objectives}

FLEXGRID directly supports the Fit-for-55 grid-flexibility target and the EU Action Plan for
Digitalising the Energy System. By keeping household consumption data on-device, the project's
differential-privacy design also anticipates the data-minimisation expectations that will apply
to any Energy Data Space use case.

\section{Measures to Maximise Impact}

\begin{minipage}{0.48\textwidth}
\begin{tcolorbox}[factbox={\faIcon{globe}\ \ Dissemination}]
\footnotesize
\begin{itemize}
  \item 6 peer-reviewed publications, 2 open-access by default
  \item Reference implementation released under Apache-2.0 at project month 30
  \item Annual DSO practitioner workshop (Amsterdam, Berlin, Madrid)
  \item Contribution to CENELEC TC205 working group
\end{itemize}
\end{tcolorbox}
\end{minipage}
\hfill
\begin{minipage}{0.48\textwidth}
\begin{tcolorbox}[factbox={\faIcon{handshake}\ \ Exploitation}]
\footnotesize
\begin{itemize}
  \item Joint exploitation agreement signed by month 6 (all 6 partners)
  \item IP: sensor design (Apex/Vertex) patent-pending, filed month 14
  \item Nexa-led spin-out business case validated in WP7
  \item Two DSOs outside consortium enrolled as early-adopter observers
\end{itemize}
\end{tcolorbox}
\end{minipage}

% ============================================================
% CHAPTER 3 -- IMPLEMENTATION
% ============================================================
\chapter{Quality and Efficiency of the Implementation}

\section{Work Plan and Work Packages}

The 36-month work plan is organised into seven work packages (WPs), sequenced so that hardware
(WP4--WP5) and software (WP2--WP3) tracks converge at the month-18 integration gate ahead of
field validation (WP6).

\begin{table}[h]
\small
\renewcommand{\arraystretch}{1.25}
\begin{tabularx}{\linewidth}{@{}l X l l@{}}
\toprule
\rowcolor{eublue}
\textcolor{white}{WP} & \textcolor{white}{Title} & \textcolor{white}{Lead} & \textcolor{white}{Months} \\
\midrule
WP1 & Project Management \& Coordination & Nexa & 1--36 \\
\rowcolor{lightgray}
WP2 & User Requirements, Ethics \& Data Governance & Meridian & 1--8 \\
WP3 & Federated Learning Core \& Edge Intelligence & Synapse & 4--24 \\
\rowcolor{lightgray}
WP4 & Resilient Grid-Edge Hardware \& Photonic Sensing & Apex & 4--24 \\
WP5 & Materials \& Durability for Extreme Environments & Vertex & 6--20 \\
\rowcolor{lightgray}
WP6 & Large-Scale Demonstration \& Field Validation & Nimbus & 18--34 \\
WP7 & Dissemination, Exploitation \& Standardisation & Nexa & 1--36 \\
\bottomrule
\end{tabularx}
\caption{Work package overview.}
\end{table}

\begin{center}
\begin{tikzpicture}[x=0.39cm,y=0.9cm,font=\sffamily]
  \def\nwp{7}
  \foreach \m in {0,6,...,36}{
    \draw[gray!35] (\m,0.4) -- (\m,-\nwp-0.2);
    \node[font=\tiny,color=slate] at (\m,0.75) {M\m};
  }
  \node[left,font=\tiny\bfseries,color=slate] at (0,-0.2) {WP1};
  \draw[fill=eublue] (0,-0.3) rectangle (36,-0.9);
  \node[left,font=\tiny\bfseries,color=slate] at (0,-1.2) {WP2};
  \draw[fill=teal] (0,-1.3) rectangle (8,-1.9);
  \node[left,font=\tiny\bfseries,color=slate] at (0,-2.2) {WP3};
  \draw[fill=teal] (4,-2.3) rectangle (24,-2.9);
  \node[left,font=\tiny\bfseries,color=slate] at (0,-3.2) {WP4};
  \draw[fill=eugold!80!black] (4,-3.3) rectangle (24,-3.9);
  \node[left,font=\tiny\bfseries,color=slate] at (0,-4.2) {WP5};
  \draw[fill=eugold!80!black] (6,-4.3) rectangle (20,-4.9);
  \node[left,font=\tiny\bfseries,color=slate] at (0,-5.2) {WP6};
  \draw[fill=slate] (18,-5.3) rectangle (34,-5.9);
  \node[left,font=\tiny\bfseries,color=slate] at (0,-6.2) {WP7};
  \draw[fill=eublue!55] (0,-6.3) rectangle (36,-6.9);
  \node[shape=diamond,draw=slate,fill=white,inner sep=1.4pt] at (18,-5.0) {};
  \node[font=\tiny,color=slate,align=left] at (18,-5.55) [xshift=1.05cm] {\tiny MS4 Integration gate};
\end{tikzpicture}
\end{center}

\subsection{Management Structure}

\begin{center}
\begin{tikzpicture}[
  node distance=7mm and 6mm,
  gbox/.style={draw=eublue,thick,rounded corners=2pt,fill=lightgray,align=center,
    font=\small\sffamily,minimum width=3.1cm,minimum height=0.95cm},
  pbox/.style={gbox,fill=white,draw=teal,minimum width=2.05cm,font=\scriptsize\sffamily}
]
\node[gbox] (pc) {Project Coordinator\\ \footnotesize Dr.\ Elena Marsh (Nexa)};
\node[gbox,below=9mm of pc] (pmo) {Project Management Office};
\node[pbox,below=9mm of pmo,xshift=-6.15cm] (p1) {Nexa\\ \footnotesize WP1, WP7};
\node[pbox,right=3.5mm of p1] (p2) {Synapse\\ \footnotesize WP3};
\node[pbox,right=3.5mm of p2] (p3) {Vertex\\ \footnotesize WP5};
\node[pbox,right=3.5mm of p3] (p4) {Nimbus\\ \footnotesize WP6};
\node[pbox,right=3.5mm of p4] (p5) {Apex\\ \footnotesize WP4};
\node[pbox,right=3.5mm of p5] (p6) {Meridian\\ \footnotesize WP2};
\draw[-{Latex[length=2mm]},thick,slate] (pc) -- (pmo);
\foreach \n in {p1,p2,p3,p4,p5,p6}{
  \draw[-{Latex[length=2mm]},thick,slate] (pmo.south) -- (\n.north);
}
\end{tikzpicture}
\end{center}

\subsection{Consortium Capacity}

\begin{table}[h]
\small
\renewcommand{\arraystretch}{1.3}
\begin{tabularx}{\linewidth}{@{}l X l l@{}}
\toprule
\rowcolor{eublue}
\textcolor{white}{Partner} & \textcolor{white}{Role \& Relevant Capacity} & \textcolor{white}{Country} & \textcolor{white}{PM} \\
\midrule
Nexa Dynamics B.V.\ (Coord.) & Federated systems integrator; coordinated 2 prior Horizon
projects & NL & 58 \\
\rowcolor{lightgray}
Synapse Analytics GmbH & Federated learning \& differential privacy research group, 40+
publications & DE & 66 \\
Vertex Materials Institute & Public research institute, polymer durability testing labs & IE & 34 \\
\rowcolor{lightgray}
Nimbus Grid Systems Oy & Mid-size DSO, operates 1{,}900\,km of MV feeder under trial & FI & 44 \\
Apex Photonics Lab AB & Photonic sensor SME, prior ERC proof-of-concept grantee & SE & 48 \\
\rowcolor{lightgray}
Meridian Energy Cooperative & Rural DSO cooperative, ethics \& community-engagement lead & ES & 32 \\
\bottomrule
\end{tabularx}
\caption{Consortium composition and effort (person-months).}
\end{table}

\section{Risk Management}

\begin{table}[h]
\small
\renewcommand{\arraystretch}{1.35}
\begin{tabularx}{\linewidth}{@{}X l l X@{}}
\toprule
\rowcolor{eublue}
\textcolor{white}{Risk} & \textcolor{white}{Likelihood} & \textcolor{white}{Impact} & \textcolor{white}{Mitigation} \\
\midrule
Field-site data access delayed by DSO internal approval & Medium & High &
\cellcolor{riskHigh}Parallel-track onboarding started at month 1; MoUs pre-negotiated \\
\addlinespace
Federated model convergence below target accuracy & Medium & Medium &
\cellcolor{riskMed}Fallback centralised-training benchmark maintained in WP3 \\
\addlinespace
Photonic sensor component supply-chain delay & Low & High &
\cellcolor{riskHigh}Dual-sourced components; 6-month buffer built into WP4 schedule \\
\addlinespace
Key personnel turnover at a partner site & Low & Medium &
\cellcolor{riskMed}Cross-training and documented handover protocol per WP \\
\addlinespace
Regulatory change to cross-border grid-data sharing & Medium & High &
\cellcolor{riskHigh}Legal monitoring by Meridian; architecture designed data-minimal by default \\
\bottomrule
\end{tabularx}
\caption{Top risks and mitigation measures.}
\end{table}

\section{Budget Summary}

\begin{table}[h]
\small
\renewcommand{\arraystretch}{1.3}
\begin{tabularx}{\linewidth}{@{}X l l@{}}
\toprule
\rowcolor{eublue}
\textcolor{white}{Beneficiary} & \textcolor{white}{Requested EU Contribution} & \textcolor{white}{Share} \\
\midrule
Nexa Dynamics B.V.\ (Coordinator) & \eur{1,150,000} & 19.2\% \\
\rowcolor{lightgray}
Synapse Analytics GmbH & \eur{1,340,000} & 22.4\% \\
Vertex Materials Institute & \eur{780,000} & 13.0\% \\
\rowcolor{lightgray}
Nimbus Grid Systems Oy & \eur{1,020,000} & 17.1\% \\
Apex Photonics Lab AB & \eur{900,000} & 15.0\% \\
\rowcolor{lightgray}
Meridian Energy Cooperative & \eur{790,000} & 13.2\% \\
\midrule
\textbf{Total} & \textbf{\eur{5,980,000}} & \textbf{100.0\%} \\
\bottomrule
\end{tabularx}
\caption{Requested EU contribution by beneficiary.}
\end{table}

\begin{tcolorbox}[notebox]
\small\faIcon{calendar}\ \ \textbf{Note on eligibility.} All effort and cost figures are drawn
up in accordance with the Horizon Europe Model Grant Agreement (Article 6, unit cost rules) and
have been reviewed by each beneficiary's certified financial officer prior to submission.
\end{tcolorbox}

\end{document}
